% Idioma y codificación
\usepackage[spanish, english, italian, es-tabla]{babel} %es-tabla para que se titule "Tabla"
\usepackage[utf8]{inputenc}

% Márgenes
\usepackage[a4paper,top=3cm,bottom=2.5cm,left=3cm,right=3cm]{geometry}

% Comentarios de bloque
\usepackage{verbatim}

% Paquetes de links
\usepackage[hidelinks]{hyperref}    % Permite enlaces
\usepackage{url}                    % redirecciona a la web

% Más opciones para enumeraciones
\usepackage{enumitem}

% Personalizar la portada
\usepackage{titling}


% Paquetes de tablas
\usepackage{multirow}


%------------------------------------------------------------------------

%Paquetes de figuras
\usepackage{caption}
\usepackage{subcaption} % Figuras al lado de otras
\usepackage{float}      % Poner figuras en el sitio indicado H.
\usepackage{multicol}


% Paquetes de imágenes
\usepackage{graphicx}       % Paquete para añadir imágenes
\usepackage{transparent}    % Para manejar la opacidad de las figuras

% Paquete para usar colores
\usepackage[dvipsnames]{xcolor}
% \usepackage{pagecolor}      % Para cambiar el color de la página

% Habilita tamaños de fuente mayores
\usepackage{fix-cm}

% Paquete para indices sin número
\usepackage{tocloft}
\renewcommand{\cftsecfont}{\normalfont} % Secciones sin negrita



%------------------------------------------------------------------------

% Paquetes de matemáticas
\usepackage{mathtools, amsfonts, amssymb, mathrsfs}
\usepackage[makeroom]{cancel}     % Simplificar tachando
\usepackage{polynom}    % Divisiones y Ruffini
\usepackage{units} % Para poner fracciones diagonales con \nicefrac

\usepackage{pgfplots}   %Representar funciones
\pgfplotsset{compat=1.18}  % Versión 1.18

\usepackage{tikz-cd}    % Para usar diagramas de composiciones
\usetikzlibrary{calc}   % Para usar cálculo de coordenadas en tikz
\usepackage{tikz}
\usetikzlibrary{positioning}
\usetikzlibrary{decorations.markings}

\usetikzlibrary{babel} % Para evitar conflictos con caracteres activos


%Definición de teoremas, etc.
\usepackage{amsthm}
%\swapnumbers   % Intercambia la posición del texto y de la numeración

% Definir un estilo de teorema con encabezado en negrita
\theoremstyle{plain} % Estilo base (puedes usar 'definition' o 'remark' si prefieres)

\newtheoremstyle{boldhead} % Nombre del estilo
  {3pt} % Espacio arriba
  {3pt} % Espacio abajo
  {\itshape} % Cuerpo del teorema (cursiva)
  {} % Indentación (vacío para ninguna)
  {\bfseries} % Estilo del encabezado (negrita)
  {.} % Puntuación después del encabezado
  { } % Espacio después del encabezado
  {} % Especificación adicional del encabezado (vacío)

% Aplicar el estilo
\theoremstyle{boldhead}

\newtheoremstyle{definition} % Nombre del estilo
  {3pt} % Espacio arriba
  {3pt} % Espacio abajo
  {} % Cuerpo del teorema (cursiva)
  {} % Indentación (vacío para ninguna)
  {\bfseries} % Estilo del encabezado (negrita)
  {.} % Puntuación después del encabezado
  { } % Espacio después del encabezado
  {} % Especificación adicional del encabezado (vacío)

% Definir el entorno para indicar nombre teorema en negrita
\makeatletter
\newenvironment{esp}[1]{%
  \begin{trivlist}
  \item[\hskip\labelsep{\bfseries #1.}]%
  \itshape % Cambiar el cuerpo a cursiva
}{%
  \end{trivlist}
}
\makeatother

\makeatletter
\@ifclassloaded{article}{
  \newtheorem{teo}{Teorema}[section]
}{
  \newtheorem{teo}{Teorema}[chapter]  % Se resetea en cada chapter
}
\makeatother

\newtheorem*{teo*}{Teorema}  % No se enumera el teorema
\newtheorem{coro}{Corolario}[teo]           % Se resetea en cada teorema
\newtheorem*{coro*}{Corolario}           % Se resetea en cada teorema
\newtheorem{prop}[teo]{Proposición}         % Usa el mismo contador que teorema
\newtheorem*{prop*}{Proposición}         % Usa el mismo contador que teorema
\newtheorem{lema}[teo]{Lema}                % Usa el mismo contador que teorema

\theoremstyle{remark}
\newtheorem*{observacion}{Observación}

\theoremstyle{definition}

\makeatletter
\@ifclassloaded{article}{
  \newtheorem{definicion}{Definición} [section]     % Se resetea en cada chapter
}{
  \newtheorem{definicion}{Definición} [chapter]     % Se resetea en cada chapter
}
\newtheorem*{definicion*}{Definición}               % Definición no numerada
\makeatother

\newtheorem*{nota*}{Note}               % Definición no numerada
\newtheorem*{nota}{Note}               % Definición no numerada
\newtheorem*{notacion}{Notación}
\newtheorem*{ejemplo}{Ejemplo}
\newtheorem*{ejercicio*}{Ejercicio}             % No numerado
\newtheorem{ejercicio}{Ejercicio} [section]     % Se resetea en cada section


% Modificar el formato de la numeración del teorema "ejercicio"
\renewcommand{\theejercicio}{%
  \ifnum\value{section}=0 % Si no se ha iniciado ninguna sección
    \arabic{ejercicio}% Solo mostrar el número de ejercicio
  \else
    \thesection.\arabic{ejercicio}% Mostrar número de sección y número de ejercicio
  \fi
}


% \renewcommand\qedsymbol{$\blacksquare$}         % Cambiar símbolo QED
%------------------------------------------------------------------------

% Paquetes para encabezados
\usepackage{fancyhdr}
\pagestyle{fancy}
\fancyhf{}

\newcommand{\helv}{ % Modificación tamaño de letra
\fontfamily{}\fontsize{12}{12}\selectfont}
\setlength{\headheight}{15pt} % Amplía el tamaño del índice


%\usepackage{lastpage}   % Referenciar última pag   \pageref{LastPage}
\fancyfoot[C]{\thepage}

% Paquete para cambiar el formato de las secciones

\usepackage{titlesec}

% Personalizar el formato de las secciones con un tamaño de fuente exacto
\titleformat{\section}
  {\fontsize{18pt}{20pt}\selectfont\bfseries} % Formato: 18pt de tamaño con 20pt de interlineado, en negrita
  {\thesection} % Mostrar el número de sección
  {1em} % Espacio entre el número y el título
  {} % Sin código adicional
  
 % Personalizar el formato de subsubsection
\titleformat{\subsubsection}
  {\large\bfseries} % Tamaño y estilo de la fuente (small y negrita)
  {\thesubsubsection} % Número de la subsección
  {1em} % Espacio entre el número y el título
  {} % Código adicional (vacío en este caso)
  
 
\titlespacing{\subsection}{0pt}{0.5cm}{0.5cm}
\titlespacing{\subsubsection}{0pt}{0.5cm}{0.5cm}

%------------------------------------------------------------------------

% Conseguir que no ponga "Capítulo 1". Sino solo "1."
\makeatletter
\@ifclassloaded{book}{
  \renewcommand{\chaptermark}[1]{\markboth{\thechapter.\ #1}{}} % En el encabezado
    
  \renewcommand{\@makechapterhead}[1]{%
  \vspace*{50\p@}%
  {\parindent \z@ \raggedright \normalfont
    \ifnum \c@secnumdepth >\m@ne
      \huge\bfseries \thechapter.\hspace{1em}\ignorespaces
    \fi
    \interlinepenalty\@M
    \Huge \bfseries #1\par\nobreak
    \vskip 40\p@
  }}
}
\makeatother

%------------------------------------------------------------------------
% Paquetes de cógido
\usepackage{minted}
\renewcommand\listingscaption{Code}

\usepackage{fancyvrb}
% Personaliza el tamaño de los números de línea
\renewcommand{\theFancyVerbLine}{\small\arabic{FancyVerbLine}}

% Estilo para C++
\newminted{cpp}{
    frame=lines,
    framesep=2mm,
    baselinestretch=1.2,
    linenos,
    escapeinside=||
}


\usepackage{listings} % Para incluir código desde un archivo

\renewcommand\lstlistingname{Code}
\renewcommand\lstlistlistingname{Índice de Códigos Fuente}

% Definir colores
\definecolor{vscodepurple}{rgb}{0.5,0,0.5}
\definecolor{vscodeblue}{rgb}{0,0,0.8}
\definecolor{vscodegreen}{rgb}{0,0.5,0}
\definecolor{vscodegray}{rgb}{0.5,0.5,0.5}
\definecolor{vscodebackground}{rgb}{0.97,0.97,0.97}
\definecolor{vscodelightgray}{rgb}{0.9,0.9,0.9}

% Configuración para el estilo de C similar a VSCode
\lstdefinestyle{vscode_C}{
  backgroundcolor=\color{vscodebackground},
  commentstyle=\color{vscodegreen},
  keywordstyle=\color{vscodeblue},
  numberstyle=\tiny\color{vscodegray},
  stringstyle=\color{vscodepurple},
  basicstyle=\scriptsize\ttfamily,
  breakatwhitespace=false,
  breaklines=true,
  captionpos=b,
  keepspaces=true,
  numbers=left,
  numbersep=5pt,
  showspaces=false,
  showstringspaces=false,
  showtabs=false,
  tabsize=2,
  frame=tb,
  framerule=0pt,
  aboveskip=10pt,
  belowskip=10pt,
  xleftmargin=10pt,
  xrightmargin=10pt,
  framexleftmargin=10pt,
  framexrightmargin=10pt,
  framesep=0pt,
  rulecolor=\color{vscodelightgray},
  backgroundcolor=\color{vscodebackground},
}

%------------------------------------------------------------------------

\usepackage[mathscr]{euscript}
\usepackage{bbm}

% Comandos para funciones concretas
\newcommand{\Log}{\text{Log}}
\newcommand{\Arg}{\text{Arg}}
\newcommand{\Arctg}{\text{Arctg}}

% Comandos definidos
\newcommand{\scr}[1]{\mathscr{#1}}
\newcommand{\bb}[1]{\mathbb{#1}}

\newcommand{\cc}[1]{\mathcal{#1}}
\newcommand{\tbf}[1]{\textbf{#1}}
\newcommand{\tnm}[1]{\textnormal{#1}}
\newcommand{\tbr}[1]{\textcolor{red}{\textbf{#1}}}

% Abreviatura para conjunto típicos
\newcommand{\E}{\mathbb{E}} % Esperanza

\newcommand{\N}{\mathbb{N}} % Conjunto de los números naturales
\newcommand{\Z}{\mathbb{Z}}
\newcommand{\Q}{\mathbb{Q}} % Conjunto de los números racionales
\newcommand{\R}{\mathbb{R}} % Conjunto de los números reales
\newcommand{\C}{\mathbb{C}} % Conjunto de los números complejos
\newcommand{\K}{\mathbb{F}} % Cuerpo cualquiera
\newcommand{\F}{\mathbb{F}} % Cuerpo cualquiera

% Abreviatura para conjuntos y funciones
\DeclareMathOperator{\GL}{GL}
\DeclareMathOperator{\SL}{SL}
\DeclareMathOperator{\Perm}{Perm}
\DeclareMathOperator{\Apl}{Apl}
\DeclareMathOperator{\ord}{ord}
\DeclareMathOperator{\Id}{Id}
\DeclareMathOperator{\Aut}{Aut}

% I prefer the slanted \leq
\let\oldleq\leq % save them in case they're every wanted
\let\oldgeq\geq
\renewcommand{\leq}{\leqslant}
\renewcommand{\geq}{\geqslant}

% Si y solo si
\newcommand{\sii}{\iff}
\newcommand{\ent}{\Longrightarrow}

% MCD y MCM
\DeclareMathOperator{\mcd}{mcd}
\DeclareMathOperator{\mcm}{mcm}

% Signo
\DeclareMathOperator{\sgn}{sgn}

% Letras griegas
\newcommand{\eps}{\epsilon}
\newcommand{\veps}{\varepsilon}
\newcommand{\lm}{\lambda}

\newcommand{\ol}{\overline}
\newcommand{\ul}{\underline}
\newcommand{\wt}{\widetilde}
\newcommand{\wh}{\widehat}

\let\oldvec\vec
\renewcommand{\vec}{\overrightarrow}

% Derivadas parciales
\newcommand{\del}[2]{\frac{\partial #1}{\partial #2}}
\newcommand{\Del}[3]{\frac{\partial^{#1} #2}{\partial^{#1} #3}}
\newcommand{\deld}[2]{\dfrac{\partial #1}{\partial #2}}
\newcommand{\Deld}[3]{\dfrac{\partial^{#1} #2}{\partial^{#1} #3}}

% Declarar aplicaciones
% 1. Nombre aplicación
% 2. Dominio
% 3. Codominio
% 4. Variable
% 5. Imagen de la variable
\newcommand{\Func}[5]{
    \begin{equation*}
        \begin{array}{rrll}
            #1:& #2 & \longrightarrow & #3\\
               & #4 & \longmapsto & #5
        \end{array}
    \end{equation*}
}

%------------------------------------------------------------------------

% Para mostrar código estético
\usepackage{csquotes}
\usepackage{xcolor} % Para definir y usar colores
\usepackage{listings} % Para mostrar código fuente
\lstset{
    language=C++,                 % Define el lenguaje del código (C++)
    basicstyle=\ttfamily\small,   % Estilo de fuente básico
    keywordstyle=\color{blue},    % Estilo de las palabras clave
    stringstyle=\color{red},      % Estilo de las cadenas de texto
    commentstyle=\color{green},   % Estilo de los comentarios
    morecomment=[l][\color{magenta}]{\#},  % Estilo para preprocesador
    numbers=left,                 % Muestra números de línea en el margen izquierdo
    numberstyle=\tiny\color{gray},% Estilo de los números de línea
    stepnumber=1,                 % Intervalo de los números de línea; cada línea
    numbersep=10pt,               % Cuánto separar los números de línea del código
    backgroundcolor=\color{white},% Color de fondo
    showspaces=false,             % Muestra espacios con subrayados; overrides 'showstringspaces'
    showstringspaces=false,       % Subraya solo los espacios en cadenas
    showtabs=false,               % Muestra tabulaciones con subrayados
    frame=single,                 % Añade un marco alrededor del código
    tabsize=1,                    % Configura tabulaciones de 2 espacios
    captionpos=b,                 % Posición de la leyenda del código (bottom)
    breaklines=true,              % Configura el quiebre automático de líneas
    breakatwhitespace=false,      % Configura el quiebre de líneas solo en espacios en blanco
    escapeinside={\%*}{*)},       % Si quieres incorporar LaTeX dentro de tu código
}

% Definición de estilo específico para JavaScript (ES6+)
\lstdefinelanguage{JavaScript}{
  keywords={break, case, catch, continue, debugger, default, delete, do, else, false, finally, for, function, if, in, instanceof, new, null, return, switch, this, throw, true, try, typeof, var, void, while, with, as, close, connect, open, send, expression, with, set, get, export, import, from, const, let},
  morekeywords=[2]{class, extends, super}, % Clases modernas
  ndkeywords={class, export, boolean, throw, implements, import, this},
  sensitive=true,
  morecomment=[l]{//},
  morecomment=[s]{/*}{*/},
  morestring=[b]',
  morestring=[b]"
}

\lstdefinestyle{vscode_JS}{
    language=JavaScript,
    basicstyle=\ttfamily\small,   % Estilo de fuente básico
    keywordstyle=\color{blue},    % Estilo de las palabras clave
    stringstyle=\color{purple},   % Estilo de las cadenas de texto (puedes usar red si prefieres)
    commentstyle=\color{green!60!black}, % Estilo de los comentarios (un verde más visible que el puro)
    morecomment=[l][\color{magenta}]{\#},
    numbers=left,                 
    numberstyle=\tiny\color{gray},
    stepnumber=1,                 
    numbersep=10pt,               
    backgroundcolor=\color{white},
    showspaces=false,             
    showstringspaces=false,       
    showtabs=false,               
    frame=single,                 
    tabsize=2,                    % Cambiado a 2 para que JavaScript no quede muy pegado
    captionpos=b,                 
    breaklines=true,              
    breakatwhitespace=false,      
    escapeinside={\%*}{*)},       
}